\section{Descomposici\'on por ra\'iz: frecuentes y poco frecuentes}
\whenuse{Objetos con frecuencias muy distintas donde procesar uno cuesta $O(f)$
o $O(n/f)$; separar pesados y livianos balancea ambas cotas.}

Usar cuando objetos con frecuencia $f$ cuestan $O(f)$ o $O(n/f)$ por consulta.
Elegir un umbral $B\approx\sqrt n$: procesar directamente los objetos con
$f<B$ y precomputar respuestas para los que tienen $f\ge B$. Hay como m\'aximo
$n/B$ objetos frecuentes.

\complexity{t\'ipicamente preproceso $O(n\sqrt n)$, consulta $O(\sqrt n)$ y
memoria entre $O(n)$ y $O(n\sqrt n)$, seg\'un lo precomputado.}

\begin{lstlisting}[style=cpp]
int sqrt_threshold(int total_items) {
    return max(1, (int)sqrt((long double)total_items) + 1);
}

vector<int> frequent_values(const vector<int> &frequency) {
    int total = accumulate(frequency.begin(), frequency.end(), 0);
    int block = sqrt_threshold(total);
    vector<int> result;
    for (int value = 0; value < (int)frequency.size(); ++value)
        if (frequency[value] >= block) result.push_back(value);
    return result;
}
\end{lstlisting}

\warning{No existe una plantilla universal: escribir primero las dos cotas y
verificar que ambas quedan balanceadas en el umbral elegido.}
